\chapter{Introducción}

\section{Sobre las notas}
Estas notas explican recordatorios y cosas basicas en programación competitiva. 
Fueron escritas con la idea de usarla como notebook en competencias. 



\section{Cosas de C++}

\subsection{Tamaño de los tipos de datos}
\begin{itemize}
	\item \textbf{Tamaño de un long long}: Es un número de 64 bits, \( 2^{64} \) o aproximadamente \( 10^{19} \).
\end{itemize}

\subsection{Lectura de líneas con espacios en C++}

En C++, para leer una línea completa que puede contener espacios (como una oración o una entrada con varios números separados), se utiliza:

\begin{lstlisting}
	getline(std::cin, s);
\end{lstlisting}

Sin embargo, si antes se usó 

\begin{lstlisting}[language=C++]
	std::cin >> variable;
\end{lstlisting}

entonces \texttt{getline()} puede leer una línea vacía. Para evitarlo, se debe limpiar el buffer con:

\begin{lstlisting}[language=C++]
	std::cin.ignore();
\end{lstlisting}



\begin{lstlisting}[language=C++]
	cin.ignore();
\end{lstlisting}



\subsection{Conversiones }

\subsubsection{String a Numero y de Numero a String }

\begin{itemize}
	\item \textbf{De string a int:}
	\texttt{int x = stoi("123");}

	\item \textbf{De string a long long:}
	\texttt{ll x = stoll("123456789");}

	\item \textbf{De string a double:}
	\texttt{double x = stod("3.14");}

	\item \textbf{De string a float:}
	\texttt{float x = stof("3.14");}

	\item \textbf{De string a long double:}
	\texttt{long double x = stold("3.14");}

	\item \textbf{De numero a string:}
	\texttt{string s = to\_string(x);}
\end{itemize}

\subsubsection{Char a Numero y de Numero a Char}

\paragraph{De char numeral a número}
Si el \texttt{char} representa un dígito (\texttt{'0'} a \texttt{'9'}):

\begin{lstlisting}[language=C++]
	int x = c - '0';
\end{lstlisting}

Ejemplo:
\begin{lstlisting}[language=C++]
	char c = '7';
	int x = c - '0'; // 7
\end{lstlisting}

\paragraph{De número a char numeral}

\begin{lstlisting}[language=C++]
	char c = x + '0';
\end{lstlisting}

Ejemplo:
\begin{lstlisting}[language=C++]
	int x = 7;
	char c = x + '0'; // '7'
\end{lstlisting}

\subsubsection{Char a Posición Alfabética y Viceversa}

\paragraph{Minúscula a número ($a \to 1$, $b \to 2$, \dots)}

\[
\text{posición}(c) = c - 'a' + 1
\]

\begin{lstlisting}[language=C++]
	int pos = c - 'a' + 1;
\end{lstlisting}

\paragraph{Número a minúscula ($1 \to a$, $2 \to b$, \dots)}

\begin{lstlisting}[language=C++]
	char c = pos - 1 + 'a';
\end{lstlisting}

\paragraph{Mayúscula a número ($A \to 1$, $B \to 2$, \dots)}

\[
\text{posición}(C) = C - 'A' + 1
\]

\begin{lstlisting}[language=C++]
	int pos = c - 'A' + 1;
\end{lstlisting}

\paragraph{Número a mayúscula ($1 \to A$, $2 \to B$, \dots)}

\begin{lstlisting}[language=C++]
	char c = pos - 1 + 'A';
\end{lstlisting}

\subsubsection{Mayúscula a Minúscula y Viceversa}

\begin{itemize}
	\item \textbf{Convertir a mayúscula}
\begin{lstlisting}[language=C++]
	char c = 'b';
	c = toupper(c); // 'B'
\end{lstlisting}

	\item \textbf{Convertir a minúscula}
\begin{lstlisting}[language=C++]
	char c = 'G';
	c = tolower(c); // 'g'
\end{lstlisting}

	\item \textbf{Verificar si es mayúscula}
\begin{lstlisting}[language=C++]
	isupper(c)
\end{lstlisting}

	\item \textbf{Verificar si es minúscula}
\begin{lstlisting}[language=C++]
	islower(c)
\end{lstlisting}
\end{itemize}

\subsubsection{Char, ASCII y Viceversa}

\paragraph{Char a código ASCII}

\begin{lstlisting}[language=C++]
	int ascii = (int)c;
\end{lstlisting}

\paragraph{ASCII a char}

\begin{lstlisting}[language=C++]
	char c = (char)x;
\end{lstlisting}

Ejemplo:
\begin{lstlisting}[language=C++]
	char c = 'A';
	int x = (int)c; // 65

	char d = (char)65; // 'A'
\end{lstlisting}

\subsubsection{String y Char}

\paragraph{Char a string}

\begin{lstlisting}[language=C++]
	char c = 'a';
	string s(1, c); // "a"
\end{lstlisting}

O también:

\begin{lstlisting}[language=C++]
	string s = string(1, c);
\end{lstlisting}

\paragraph{String de longitud 1 a char}

\begin{lstlisting}[language=C++]
	char c = s[0];
\end{lstlisting}

\subsubsection{String a Vector de Chars y Viceversa}

\paragraph{String a \texttt{vector<char>}}

\begin{lstlisting}[language=C++]
	string s = "hola";
	vector<char> v(s.begin(), s.end());
\end{lstlisting}

\paragraph{\texttt{vector<char>} a string}

\begin{lstlisting}[language=C++]
	string s(v.begin(), v.end());
\end{lstlisting}

\subsubsection{Numero Decimal a Otras Bases}

\paragraph{Decimal a binario}

\begin{lstlisting}[language=C++]
	ll x = 13;
	string s = bitset<64>(x).to_string();

	s.erase(0, s.find('1'));
\end{lstlisting}

\paragraph{Decimal a hexadecimal}

\begin{lstlisting}[language=C++]
	stringstream ss;
	ss << hex << x;
	string hexNum = ss.str();
\end{lstlisting}

\paragraph{Decimal a octal}

\begin{lstlisting}[language=C++]
	stringstream ss;
	ss << oct << x;
	string octNum = ss.str();
\end{lstlisting}

\subsubsection{String en Base N a Decimal}

\paragraph{Binario a decimal}

\begin{lstlisting}[language=C++]
	ll x = stoll(binary, nullptr, 2);
\end{lstlisting}

\paragraph{Hexadecimal a decimal}

\begin{lstlisting}[language=C++]
	ll x = stoll(hexString, nullptr, 16);
\end{lstlisting}

\paragraph{Octal a decimal}

\begin{lstlisting}[language=C++]
	ll x = stoll(octalString, nullptr, 8);
\end{lstlisting}

\subsubsection{Numeros entre Tipos}

\paragraph{int a long long}

\begin{lstlisting}[language=C++]
	ll x = (ll)a;
\end{lstlisting}

\paragraph{double a int}

\begin{lstlisting}[language=C++]
	int x = (int)d;
\end{lstlisting}

\paragraph{int a double}

\begin{lstlisting}[language=C++]
	double x = (double)a;
\end{lstlisting}

\paragraph{División real entre enteros}

\begin{lstlisting}[language=C++]
	double ans = (double)a / b;
\end{lstlisting}

\subsubsection{String a Array de Digitos}

\paragraph{String de dígitos a \texttt{vector<int>}}

\begin{lstlisting}[language=C++]
	string s = "12345";
	vector<int> digits;

	for(char c : s){
		digits.push_back(c - '0');
	}
\end{lstlisting}

\paragraph{\texttt{vector<int>} de dígitos a string}

\begin{lstlisting}[language=C++]
	string s;

	for(int d : digits){
		s += char(d + '0');
	}
\end{lstlisting}

\subsection*{Ordenar partes (subrangos) de un vector}

En C++, se puede ordenar una parte de un vector (es decir, un subrango) utilizando la función \texttt{sort} de la STL, especificando iteradores de inicio y fin.

Esto es especialmente útil cuando se necesita dividir un conjunto en mitades o grupos, por ejemplo, para aplicar diferentes criterios de ordenamiento en cada parte.

\begin{lstlisting}[language=C++]
	// Supongamos que tenemos un vector de arreglos de 3 enteros:
	vector<array<ll, 3>> v(n);
	
	// Ordenamos solo la primera mitad del vector por el segundo campo (indice 1):
	sort(v.begin(), v.begin() + n/2, [](const array<ll, 3> &a, const array<ll, 3> &b) {
		return a[1] < b[1];
	});
	
	// Ordenamos la segunda mitad tambien por el segundo campo:
	sort(v.begin() + n/2, v.end(), [](const array<ll, 3> &a, const array<ll, 3> &b) {
		return a[1] < b[1];
	});
\end{lstlisting}

\textbf{Aplicación:} esta técnica aparece frecuentemente en problemas donde hay que emparejar elementos de dos mitades, o donde diferentes partes del conjunto deben ser tratadas con diferentes estrategias.

\section{Cosas comunes en programación competitiva}


\subsection{Notación BIG O}
\begin{table}[h!]
	\centering
	\begin{tabular}{|c|c|}
		\hline
		\textbf{Tamaño del input: $N$} & \textbf{Peor complejidad aceptada} \\ \hline
		$N \leq 22 \ (\approx 2 \cdot 10^1)$ & $O(2^N)$ \\ \hline
		$N \leq 100 \ (10^2)$ & $O(N^3)$ \\ \hline
		$N \leq 1000 \ (10^3)$ & $O(N^2)$ \\ \hline
		$N \leq 1{,}000{,}000 \ (10^6)$ & $O(N \log N)$ \\ \hline
		$N \leq 1{,}000{,}000{,}000 \ (10^9)$ & $O(N),\ O(\log N),\ O(1)$ \\ \hline
	\end{tabular}
	\caption{Relación entre tamaño de input y la peor complejidad aceptada.}
\end{table}

\subsection{Resize() en c++}
La funci\'on \texttt{std::vector::resize()} ajusta el n\'umero de elementos en un vector. Su acci\'on depende del tama\~no deseado:

\begin{itemize}
	\item \textbf{Si el nuevo tama\~no es mayor que el actual:} Se a\~naden nuevos elementos al final. Opcionalmente, se puede especificar un valor de inicializaci\'on: \texttt{mi\_vector.resize(10, -1)}.
	\item \textbf{Si el nuevo tama\~no es menor que el actual:} Se eliminan los elementos del final del vector hasta alcanzar el tama\~no especificado.
\end{itemize}
\begin{lstlisting}[style=mycpp, caption=Uso eficiente de resize]

	int main() {
		int n; cin>>n;
		vector<int> numeros;
		numeros.resize(n); // Pre-asigna memoria para n elementos
		for (int i = 0; i < n; ++i) {
			cin >> numeros[i]; // Acceso directo por \'indice
		}
		return 0;
	}
\end{lstlisting}

